\documentclass[10pt,a4paper]{amsart}
\usepackage{amsmath,amssymb,amsthm,graphicx}
% Retained preamble for the Stacks Project chapter morphisms.tex (Morphisms of Schemes).
% The chapter file pinned-original-source.tex carries only \input{preamble} on its first line, so no
% preamble text lives inside the chapter itself.  The Stacks Project uses one single shared preamble for
% every chapter; the snapshot stages it verbatim in sources/stacks-chow/preamble.tex.  The block below is
% copied mechanically, line for line, from that file's lines 14-18 (the Xy-pic setup) and lines 154-217
% (the native theorem declarations together with the mathematical macros).  Nothing was added, removed or
% redefined.  The wrapper class and the amsmath/amssymb/amsthm/graphicx base packages are supplied by the
% assembler, so the class/\externaldocument part of the shared preamble is not repeated here.
% For commutative diagrams we use Xy-pic
\usepackage[all]{xy}

% We use 2cell for 2-commutative diagrams.
\xyoption{2cell}
% Theorem environments.
%
\theoremstyle{plain}
\newtheorem{theorem}[subsection]{Theorem}
\newtheorem{proposition}[subsection]{Proposition}
\newtheorem{lemma}[subsection]{Lemma}

\theoremstyle{definition}
\newtheorem{definition}[subsection]{Definition}
\newtheorem{example}[subsection]{Example}
\newtheorem{exercise}[subsection]{Exercise}
\newtheorem{situation}[subsection]{Situation}

\theoremstyle{remark}
\newtheorem{remark}[subsection]{Remark}
\newtheorem{remarks}[subsection]{Remarks}

\numberwithin{equation}{subsection}

% Macros
%
\def\lim{\mathop{\mathrm{lim}}\nolimits}
\def\colim{\mathop{\mathrm{colim}}\nolimits}
\def\Spec{\mathop{\mathrm{Spec}}}
\def\Hom{\mathop{\mathrm{Hom}}\nolimits}
\def\Ext{\mathop{\mathrm{Ext}}\nolimits}
\def\SheafHom{\mathop{\mathcal{H}\!\mathit{om}}\nolimits}
\def\SheafExt{\mathop{\mathcal{E}\!\mathit{xt}}\nolimits}
\def\Sch{\mathit{Sch}}
\def\Mor{\mathop{\mathrm{Mor}}\nolimits}
\def\Ob{\mathop{\mathrm{Ob}}\nolimits}
\def\Sh{\mathop{\mathit{Sh}}\nolimits}
\def\NL{\mathop{N\!L}\nolimits}
\def\CH{\mathop{\mathrm{CH}}\nolimits}
\def\proetale{{pro\text{-}\acute{e}tale}}
\def\etale{{\acute{e}tale}}
\def\QCoh{\mathit{QCoh}}
\def\Ker{\mathop{\mathrm{Ker}}}
\def\Im{\mathop{\mathrm{Im}}}
\def\Coker{\mathop{\mathrm{Coker}}}
\def\Coim{\mathop{\mathrm{Coim}}}

% Boxtimes
%
\DeclareMathSymbol{\boxtimes}{\mathbin}{AMSa}{"02}

%
% Macros for moduli stacks/spaces
%
\def\QCohstack{\mathcal{QC}\!\mathit{oh}}
\def\Cohstack{\mathcal{C}\!\mathit{oh}}
\def\Spacesstack{\mathcal{S}\!\mathit{paces}}
\def\Quotfunctor{\mathrm{Quot}}
\def\Hilbfunctor{\mathrm{Hilb}}
\def\Curvesstack{\mathcal{C}\!\mathit{urves}}
\def\Polarizedstack{\mathcal{P}\!\mathit{olarized}}
\def\Complexesstack{\mathcal{C}\!\mathit{omplexes}}
% \Pic is the operator that assigns to X its picard group, usage \Pic(X)
% \Picardstack_{X/B} denotes the Picard stack of X over B
% \Picardfunctor_{X/B} denotes the Picard functor of X over B
\def\Pic{\mathop{\mathrm{Pic}}\nolimits}
\def\Picardstack{\mathcal{P}\!\mathit{ic}}
\def\Picardfunctor{\mathrm{Pic}}
\def\Deformationcategory{\mathcal{D}\!\mathit{ef}}

% --- appended 2026-10-12 (Chief): the project own definitions of the reference/slogan/history
% display environments, copied mechanically from the shared Stacks preamble lines 8-12
% (\usepackage{verbatim} plus the three \newenvironment lines, which the chapter preamble keeps
% verbatim from the project preamble). Nothing added, removed or redefined: these environments
% carry no mathematical content and the project itself typesets them as comments. Without these
% three lines a counted statement whose own text uses them fails with Environment slogan undefined.
\usepackage{verbatim}
\newenvironment{reference}{\comment}{\endcomment}
\newenvironment{slogan}{\comment}{\endcomment}
\newenvironment{history}{\comment}{\endcomment}
\usepackage{cleveref}
\crefname{originallemma}{original lemma}{original lemmas}
\Crefname{originallemma}{Original lemma}{Original lemmas}
\makeatletter
\newlabel{algebra-lemma-field-finite-type-over-domain}{{00FY}{source}{}{sourcepointer-1}{}}
\newlabel{algebra-lemma-field-finite-type-over-domain@cref}{{[originallemma][0][]00FY}{[1][0][]source}{}{}{}}
\makeatother
\crefname{originallemma}{original lemma}{original lemmas}
\Crefname{originallemma}{Original lemma}{Original lemmas}
\makeatletter
\newlabel{algebra-lemma-integral-under-field}{{00GR}{source}{}{sourcepointer-2}{}}
\newlabel{algebra-lemma-integral-under-field@cref}{{[originallemma][0][]00GR}{[1][0][]source}{}{}{}}
\makeatother
\crefname{originallemma}{original lemma}{original lemmas}
\Crefname{originallemma}{Original lemma}{Original lemmas}
\makeatletter
\newlabel{lemma-locally-finite-type-characterize}{{01T2}{source}{}{sourcepointer-3}{}}
\newlabel{lemma-locally-finite-type-characterize@cref}{{[originallemma][0][]01T2}{[1][0][]source}{}{}{}}
\makeatother
\title{Four equivalent conditions for a morphism from the spectrum of a field to be of finite type}
\author{The Stacks Project Authors}
\date{stacks-project:morphisms.tex:local-snapshot-20261012; GNU FDL 1.2 or later}
\begin{document}
\maketitle
\noindent\textit{Editorial note.} All delivered mathematical text is the author's own, copied from the pinned Stacks Project chapter file pinned-original-source.tex; no statement or proof step was rewritten, reordered or removed.\par\medskip
\noindent\textit{Editorial note.} Following the multi-claim convention for this source, source\_url is the source-level URL https://stacks.math.columbia.edu/ and the source label lemma-point-finite-type is carried in claim\_key and here rather than in a per-claim tag URL.\par\medskip
\noindent\textit{Editorial note.} The preamble is the chapter's shared Stacks Project preamble, retained as sources/stacks-morphisms/retained-preamble.tex: a line-for-line excerpt of the shared Stacks preamble staged with this snapshot. No macro was invented or redefined, and this card adds no preamble text of its own.\par\medskip
\noindent\textit{Editorial note.} Statement span 2826-2840 and proof span 2842-2875 come from the halfA pool card operations/pool-cards/stacks-morphisms-halfA.json and were re-read in the pinned chapter file before this card was written.\par\medskip
\noindent\textit{Editorial note.} The proof uses the letters (1) to (4) as labels for the four conditions of the statement; they are the author's and are kept verbatim.\par\medskip
\noindent\textit{Editorial note.} 3 cross-references occur inside the counted spans; they are listed in external\_references as navigation pointers only.\par\medskip
\noindent\textit{Allowed standard prerequisites (editorial).}\begin{itemize}
\item A morphism from the spectrum of a field is quasi-compact because its source is a singleton, so finite type and locally finite type agree for it. Used at proof lines 3059-3061. Gives the equivalence of the first two conditions.
\item The algebra lemma that a finite type domain extension inside a field has a localization that is a field with finite extension of the target field, cited from the algebra chapter. Cited at proof lines 3068-3070. Produces the element whose localization turns the finite type injection into a finite field extension, giving the third condition.
\item The same-chapter characterization of locally finite type morphisms on affine charts, and the algebra lemma that the image of a finite ring map into a field is a field over which the target is finite. Cited at proof lines 3074-3075 and 3079-3080. Give the implication from the third to the second condition, and the identification of the kernel of a finite ring map into a field as a maximal ideal.
\end{itemize}\medskip
\ifdefined\hypertarget\hypertarget{sourcepointer-1}{}\fi
\noindent\textit{Source reference.} Original lemma 00FY: focused-seeds/AGP/sources/stacks-snapshot-20261011/tags.txt, source lines 588--588. Same-chapter reference used in the retained author text at source line(s) 2852. It is supplied only as a navigation pointer so that the retained text compiles; the counted item delivers the author's own statement and the author's complete proof, and no part of the referenced item is copied in or substituted for the counted argument.\par\medskip
\ifdefined\hypertarget\hypertarget{sourcepointer-2}{}\fi
\noindent\textit{Source reference.} Original lemma 00GR: focused-seeds/AGP/sources/stacks-snapshot-20261011/tags.txt, source lines 617--617. Same-chapter reference used in the retained author text at source line(s) 2869. It is supplied only as a navigation pointer so that the retained text compiles; the counted item delivers the author's own statement and the author's complete proof, and no part of the referenced item is copied in or substituted for the counted argument.\par\medskip
\ifdefined\hypertarget\hypertarget{sourcepointer-3}{}\fi
\noindent\textit{Source reference.} Original lemma 01T2: focused-seeds/AGP/sources/stacks-snapshot-20261011/tags.txt, source lines 2359--2359. Same-chapter reference used in the retained author text at source line(s) 2717, 2739, 2801, 2863, 3019, 3035, 3157, 3348, 3613, 3707, 3868, 5197, 5282, 6702, 9325, 9390, 9440, 11421, 13422, 13686. It is supplied only as a navigation pointer so that the retained text compiles; the counted item delivers the author's own statement and the author's complete proof, and no part of the referenced item is copied in or substituted for the counted argument.\par\medskip
\begin{lemma}
\label{lemma-point-finite-type}
Let $S$ be a scheme. Let $k$ be a field.
Let $f : \Spec(k) \to S$ be a morphism.
The following are equivalent:
\begin{enumerate}
\item The morphism $f$ is of finite type.
\item The morphism $f$ is locally of finite type.
\item There exists an affine open $U = \Spec(R)$ of $S$
such that $f$ corresponds to a finite ring map $R \to k$.
\item There exists an affine open $U = \Spec(R)$ of $S$
such that the image of $f$ consists of a closed point $u$ in $U$
and the field extension $k/\kappa(u)$ is finite.
\end{enumerate}
\end{lemma}

\begin{proof}
The equivalence of (1) and (2) is obvious as $\Spec(k)$
is a singleton and hence any morphism from it is quasi-compact.

\medskip\noindent
Suppose $f$ is locally of finite type. Choose any affine open
$\Spec(R) = U \subset S$ such that the image of $f$
is contained in $U$, and the ring map $R \to k$
is of finite type. Let $\mathfrak p \subset R$ be the kernel.
Then $R/\mathfrak p \subset k$ is of finite type. By
Algebra, Lemma \ref{algebra-lemma-field-finite-type-over-domain}
there exist a $\overline{f} \in R/\mathfrak p$ such that
$(R/\mathfrak p)_{\overline{f}}$ is a field and
$(R/\mathfrak p)_{\overline{f}} \to k$ is a finite field
extension. If $f \in R$ is a lift of $\overline{f}$, then
we see that $k$ is a finite $R_f$-module. Thus (2) $\Rightarrow$ (3).

\medskip\noindent
Suppose that $\Spec(R) = U \subset S$ is an affine open
such that $f$ corresponds to a finite ring map $R \to k$.
Then $f$ is locally of finite type
by Lemma \ref{lemma-locally-finite-type-characterize}.
Thus (3) $\Rightarrow$ (2).

\medskip\noindent
Suppose $R \to k$ is finite. The image of $R \to k$ is a field
over which $k$ is finite by
Algebra, Lemma \ref{algebra-lemma-integral-under-field}.
Hence the kernel of $R \to k$ is a maximal ideal.
Thus (3) $\Rightarrow$ (4).

\medskip\noindent
The implication (4) $\Rightarrow$ (3) is immediate.
\end{proof}

\end{document}
